% Fragmento integrable. Preambulo anfitrion: amsmath, amssymb, longtable, hyperref.
\section{Entrada transversal y escalado local}
\label{hmtfg:fragmento:entrada-transversal-y-escalado-local}

\subsection{Procedencia y resultado}
\label{hmtfg:escalado:1}

APP evalúa suma y producto sobre las dos hojas de la retícula, conservando cociente y residuo. TRIT determina régimen y orientación. El transporte TPK selecciona, transporta y actualiza el estado con sus registros de acarreo, frontera y memoria. Sus lectores actúan sobre la misma estructura discreta del continuo. La prolongación \texttt{w6→w12→w18→w24→w30→R36→G9} conserva esta procedencia y distingue retorno de fase de incremento de memoria.

El lector crítico de la rueda dodecafásica, construido en los antecedentes operatorios y transportado por el refinamiento compatible, produce
\[
\phi_j=\frac{(3j-1)\pi_{\mathrm{HMT}}}{18},\qquad
w_j=\frac1{12},\qquad
k_j=\sqrt{\frac2{\sum_{\ell=1}^{12}w_\ell\phi_\ell^2}}\,\phi_j
=\frac{2(3j-1)}{\sqrt{899}}.
\]
La coordenada \(\pi_{\mathrm{HMT}}\) conserva su origen APP–TRIT–TPK; su factor común se cancela en este segundo momento. El lector resultante es
\begin{equation}
u_0(x)=\frac1{12}\sum_{j=1}^{12}(1-\cos k_jx),\qquad
f_\lambda(x)=1-\lambda u_0(x).
\label{hmtfg:escalado:eq:1}
\end{equation}
El sector focal es el uniforme, \(\eta=0\). La ampliación antecedente conserva por separado la familia ponderada \(|\eta|\le1/40\). Las nuevas cotas numéricas corresponden a \eqref{hmtfg:escalado:eq:1}.

La familia es par, entera, con \(f_\lambda(0)=1\), crítico cuadrático y segundo momento \(\sum_jw_jk_j^2=2\). El parámetro \(\lambda\) recorre una familia posterior al generador HMT. La constante objetivo de Feigenbaum permanece fuera de las entradas de \eqref{hmtfg:escalado:eq:1}.

\textbf{Resultado principal.} Existe un único
\begin{equation}
\lambda_*\in[1.874038,1.874039]
\label{hmtfg:escalado:eq:2}
\end{equation}
para el que el cuarto retorno normalizado de \eqref{hmtfg:escalado:eq:1} pertenece a la hoja estable local del punto fijo real de duplicación. El cruce de la familia con esa hoja es transversal, con separación cuantificada. Para todo \(n\ge0\),
\begin{equation}
\boxed{\|R^{4+n}f_{\lambda_*}-g\|_{\mathcal A}
<0.001666\,(0.83996)^n.}
\label{hmtfg:escalado:eq:3}
\end{equation}
La cascada superestable local asociada a este cruce posee parámetros \(\mu_j\) y constantes \(C\ne0\), \(0<\theta<1\), tales que
\begin{equation}
\mu_j-\lambda_*=C\delta_{\mathrm F}^{-j}\bigl[1+O(\theta^j)\bigr],
\qquad
\frac{\mu_{j-1}-\mu_{j-2}}{\mu_j-\mu_{j-1}}\longrightarrow\delta_{\mathrm F}.
\label{hmtfg:escalado:eq:4}
\end{equation}
Aquí \(\delta_{\mathrm F}>1\) es el autovalor inestable del operador de duplicación. La selección de \(\mu_j\) se realiza mediante los retornos locales de sección~\ref{hmtfg:escalado:8}. La identificación con el selector global se demuestra en secciones~\ref{hmtfg:escalado:15}--\ref{hmtfg:escalado:17}: todos los términos globales conservan la combinatoria y los centros coinciden para todo nivel suficientemente alto al igualar sus períodos.


\subsection{Carta analítica y retorno efectivo}
\label{hmtfg:escalado:2}

Se emplea la carta de funciones pares
\[
s=x^2,\qquad z=\frac{s-1}{5/2},\qquad
f(x)=1-x^2V(z),\qquad V(z)=\sum_{j\ge0}v_jz^j.
\]
Las coordenadas del espacio de Banach son
\begin{equation}
u=10v_0,\qquad \nu=(v_1,v_2,\ldots),\qquad
\|(\Delta u,\Delta\nu)\|_{\mathcal A}
=|\Delta u|+\sum_{j\ge1}|\Delta v_j|.
\label{hmtfg:escalado:eq:5}
\end{equation}
En particular, el coeficiente constante de \(V\) tiene peso diez. Esta normalización coincide con la carta \(u/10\) de Hertling–Spandl, sección 2 de la fuente.

Sea \(a=-f(1)=v_0-1\). Para los retornos reales admisibles,
\begin{equation}
(Rf)(x)=-\frac{f(f(-ax))}{a}.
\label{hmtfg:escalado:eq:6}
\end{equation}
El programa evalúa una factorización exacta de \eqref{hmtfg:escalado:eq:6}. Definiendo
\[
S=\frac{a^2s-1}{5/2},\qquad b=V(S),\qquad
h=1-a^2sb,\qquad t=\frac{h^2-1}{5/2},
\]
\[
\operatorname{shift}V(z)=\sum_{j\ge0}v_{j+1}z^j,
\]
resulta
\begin{equation}
\boxed{V_R=a\,b\,(h+1)\left[V(t)+\frac25\operatorname{shift}V(t)\right].}
\label{hmtfg:escalado:eq:7}
\end{equation}
Para comprobarla, se usa \(v_0=1+a\),
\(h^2-1=-a^2sb(h+1)\) y
\(V(t)-v_0=t\operatorname{shift}V(t)\). Así,
\(a+1-h^2V(t)=a^2sb(h+1)[V(t)+(2/5)\operatorname{shift}V(t)]\),
que es precisamente \(asV_R\).

Los momentos iniciales son racionales:
\begin{equation}
\frac1{12}\sum_jk_j^{2m}
=\frac1{12}\left(\frac4{899}\right)^m
\sum_{j=1}^{12}(3j-1)^{2m}.
\label{hmtfg:escalado:eq:8}
\end{equation}
Por ello, tanto \eqref{hmtfg:escalado:eq:7} como la construcción inicial utilizan intervalos racionales; las funciones trigonométricas iniciales se representan por sus series con resto exterior.


\subsection{Cotas externas utilizadas como reconocimiento posterior}
\label{hmtfg:escalado:3}

Una vez construida \eqref{hmtfg:escalado:eq:1}, se utiliza el teorema de hiperbolicidad de Lanford como certificado analítico del operador \eqref{hmtfg:escalado:eq:6}. Sean \(g_0\) el polinomio central publicado y \(g\) el punto fijo. Las estimaciones utilizadas son
\[
\|g-g_0\|_{\mathcal A}<\varepsilon_0=4\,10^{-5},
\]
y, en la bola \(\|f-g_0\|_{\mathcal A}<0.01\), los bloques de \(DR\), escritos \(R=(U,V)\), satisfacen
\begin{equation}
U_u\ge a_0=4.521,\quad
\|U_\nu\|\le b_0=0.560,\quad
\|V_u\|\le c_0=0.756,\quad
\|V_\nu\|\le d_0=0.719.
\label{hmtfg:escalado:eq:9}
\end{equation}
El punto fijo posee una única dirección inestable, con autovalor real positivo \(\delta_{\mathrm F}>1\); el resto del espectro queda dentro del disco unidad. Las cifras de \eqref{hmtfg:escalado:eq:9} proceden de las estimaciones publicadas, no de una nueva reproducción de la prueba de Lanford en este expediente. Los coeficientes de \(g_0\) están incorporados explícitamente en el polinomio central de referencia del programa, conforme a la tabla 2 de Hertling–Spandl.

Fuentes: \href{https://repo-archives.ihes.fr/A_GARDER/20180409/Test/I_Prepublications/LANFORD/1968-2013/P_81_17/P_81_17_web.pdf}{Lanford, estimaciones de hiperbolicidad}; \href{https://arxiv.org/pdf/1410.3277}{Hertling–Spandl, carta analítica y tabla 2}.


\subsection{Certificado racional de entrada y dirección}
\label{hmtfg:escalado:4}

El certificado de entrada a la renormalización divide \eqref{hmtfg:escalado:eq:2} en dieciséis intervalos racionales iguales. Propaga cuatro retornos de \eqref{hmtfg:escalado:eq:7} y la derivada exacta respecto a \(\lambda\), usando diferenciación de la composición y redondeo exterior a cien posiciones decimales.

Cada serie consta de un polinomio de grado 32 con coeficientes intervalares y un error analítico arbitrario \(E\), \(\|E\|_1\le\epsilon\). Este error puede modificar también coeficientes bajos. Para \(\|q\|_1<1\),
\begin{equation}
\|E\circ q\|_1\le\epsilon,\qquad
\|E'\circ q\|_1\le\frac{\epsilon}{(1-\|q\|_1)^2},\qquad
\|\operatorname{shift}E\|_1\le\epsilon.
\label{hmtfg:escalado:eq:10}
\end{equation}
Los productos incluyen los términos polinomio–error y error–error. La cola inicial se acota por su primer término absoluto y una razón geométrica, utilizando \(\|1+(5/2)z\|_1=7/2\). Cada argumento compuesto permanece estrictamente dentro de la bola unidad de esta álgebra.

Para \(\Gamma(\lambda)=R^4f_\lambda=(u_\lambda,\nu_\lambda)\), el certificado da las siguientes cotas conservadoras, uniformes en \eqref{hmtfg:escalado:eq:2}:

\begingroup\small
\begin{longtable}{p{0.465\linewidth}p{0.465\linewidth}}
\hline
Cantidad & Cota certificada exterior \\
\hline
\(\|\Gamma(\lambda)-g_0\|_{\mathcal A}\) & \(<0.004993128\) \\
\(\|\nu_\lambda-\nu_{g_0}\|_1\) & \(<0.001396077\) \\
\(u'_\lambda\) & \(>3821.289115\) \\
\(\|\nu'_\lambda\|_1\) & \(<551.757000\) \\
\(\|\nu'_\lambda\|_1/u'_\lambda\) & \(<0.144386<1/4\) \\
Error analítico de la función & \(<8.482\,10^{-19}\) \\
Error analítico de la tangente & \(<4.569\,10^{-15}\) \\
\hline\end{longtable}
\endgroup

Los decimales completos, los dieciséis subintervalos y cada retorno intermedio se conservan en el certificado racional de entrada.

El cono \(\|\dot\nu\|_1\le|\dot u|/4\) es invariante bajo \eqref{hmtfg:escalado:eq:9}:
\[
|\dot u_{\rm nuevo}|\ge(4.521-0.560/4)|\dot u|
=4.381|\dot u|,
\]
\begin{equation}
\|\dot\nu_{\rm nueva}\|_1\le(0.756+0.719/4)|\dot u|
=0.93575|\dot u|<\tfrac14(4.381)|\dot u|.
\label{hmtfg:escalado:eq:11}
\end{equation}
La dirección paramétrica certificada entra, por tanto, en un cono uniformemente expansivo del operador.


\subsection{Construcción cuantitativa de la hoja estable}
\label{hmtfg:escalado:5}

Traslademos \(g\) al origen y escribamos las coordenadas como \((x,y)=(u-u_g,\nu-\nu_g)\). Fijemos
\begin{equation}
L=0.16,\qquad r=0.004,\qquad
A=a_0-Lc_0=4.40004,\qquad q=d_0+c_0L=0.83996.
\label{hmtfg:escalado:eq:12}
\end{equation}
El cilindro \(|x|\le Lr\), \(\|y\|_1\le r\) está dentro de la bola de \eqref{hmtfg:escalado:eq:9}, puesto que
\((1+L)r+\varepsilon_0=0.00468<0.01\).

Considérese el espacio completo de funciones \(\sigma:B_r\to\mathbb R\) con \(\sigma(0)=0\) y constante de Lipschitz a lo sumo \(L\), dotado de la distancia uniforme. Para cada \(y\), se define \((\mathcal G\sigma)(y)\) como la raíz en \([-L\|y\|_1,L\|y\|_1]\) de
\begin{equation}
F_\sigma(x,y)=U(x,y)-\sigma(V(x,y)).
\label{hmtfg:escalado:eq:13}
\end{equation}
En ese intervalo, \(\|V(x,y)\|_1\le q\|y\|_1<r\). Al incrementar \(x\), \eqref{hmtfg:escalado:eq:9} muestra que \(F_\sigma\) crece al menos con pendiente \(A\). En sus extremos,
\[
F_\sigma(L\|y\|_1,y)
\ge[AL-b_0-Ld_0]\|y\|_1,
\]
y la desigualdad opuesta vale en el extremo negativo. El margen es
\begin{equation}
AL-b_0-Ld_0=0.0289664>0.
\label{hmtfg:escalado:eq:14}
\end{equation}
Se obtiene una raíz única. Comparando \eqref{hmtfg:escalado:eq:13} para dos valores de \(y\),
\[
\operatorname{Lip}(\mathcal G\sigma)
\le\frac{b_0+Ld_0}{A}
=\frac{0.67504}{4.40004}<0.16.
\]
Comparando dos gráficos,
\begin{equation}
\|\mathcal G\sigma-\mathcal G\widetilde\sigma\|_\infty
\le A^{-1}\|\sigma-\widetilde\sigma\|_\infty,
\qquad A^{-1}<0.228.
\label{hmtfg:escalado:eq:15}
\end{equation}
El teorema de contracción produce un único gráfico fijo \(x=\sigma_*(y)\). Sobre él,
\begin{equation}
\|y_n\|_1\le q^n\|y_0\|_1,
\qquad
\|(x_n,y_n)\|_{\mathcal A}\le(1+L)q^n\|y_0\|_1.
\label{hmtfg:escalado:eq:16}
\end{equation}
Además, la distancia vertical \(d(x,y)=x-\sigma_*(y)\) cumple
\[
|d(R(x,y))|\ge A|d(x,y)|
\]
mientras la órbita permanece en el cilindro. Toda órbita que permanece allí indefinidamente pertenece al gráfico. Este identifica exactamente la hoja estable local y proporciona la estimación \eqref{hmtfg:escalado:eq:16}. Su regularidad local coincide con la variedad estable analítica del punto fijo hiperbólico.


\subsection{Existencia, unicidad y transversalidad del parámetro}
\label{hmtfg:escalado:6}

Defínase, para \(\lambda\) en \eqref{hmtfg:escalado:eq:2},
\[
H(\lambda)=u_\lambda-u_g-
\sigma_*(\nu_\lambda-\nu_g).
\]
El certificado asegura \(\|\nu_\lambda-\nu_g\|_1<0.001436077<r\). Para \(\lambda_2>\lambda_1\),
\begin{equation}
H(\lambda_2)-H(\lambda_1)
\ge\int_{\lambda_1}^{\lambda_2}
\bigl[u'_t-L\|\nu'_t\|_1\bigr]dt
>3733.007995\,(\lambda_2-\lambda_1).
\label{hmtfg:escalado:eq:17}
\end{equation}
Para decidir los signos extremos se utiliza \(\|g-g_0\|_{\mathcal A}<\varepsilon_0\) y \(|\sigma_*(y)|\le L\|y\|_1\). La evaluación puntual racional produce
\begin{equation}
H(1.874038)<-0.0011856110945,
\qquad H(1.874039)>0.0021866053399.
\label{hmtfg:escalado:eq:18}
\end{equation}
Por continuidad, existe una raíz; \eqref{hmtfg:escalado:eq:17} prueba su unicidad y la transversalidad del cruce. En esa raíz, \eqref{hmtfg:escalado:eq:16} y
\((1+L)(0.001396077+0.00004)<0.001666\)
demuestran \eqref{hmtfg:escalado:eq:3}.

La cota es uniforme en el número de retornos: convierte el avance de profundidad en un error operatorial geométrico explícito.


\subsection{Admisibilidad real de todos los retornos}
\label{hmtfg:escalado:7}

Los cuatro retornos iniciales se acompañan, en cada subintervalo racional, de las desigualdades
\begin{equation}
0<a=-f(1)<1,\qquad b=f(a)>a,\qquad f(b)<a.
\label{hmtfg:escalado:eq:19}
\end{equation}
El programa guarda \(a,b,f(b)\) antes de cada aplicación de \eqref{hmtfg:escalado:eq:7}. La unimodalidad, la paridad y el crítico cuadrático se conservan bajo estos retornos normalizados.

Para las iteraciones siguientes basta una verificación uniforme en la bola de \eqref{hmtfg:escalado:eq:9}. Si \(E=V-V_{g_0}\), la norma \eqref{hmtfg:escalado:eq:5} da, para \(-0.4\le z\le0\),
\[
|E(z)|\le0.004,\qquad |E'(z)|\le0.01.
\]
La segunda desigualdad usa \(\sup_{j\ge1}j(0.4)^{j-1}=1\). La evaluación racional de los coeficientes centrales proporciona
\begin{equation}
\frac65<V(z)<\frac85,\qquad |V'(z)|<\frac12,
\qquad 0.39<a<0.41.
\label{hmtfg:escalado:eq:20}
\end{equation}
En consecuencia,
\[
f'(x)=-2x\left[V(z)+\frac25x^2V'(z)\right],
\]
y el corchete supera uno para \(0\le x\le1\). El mapa es unimodal y preserva \([-1,1]\). Además,
\[
b=f(a)>1-(0.41)^2\frac85=\frac{4569}{6250}>a,
\]
\begin{equation}
f(b)<1-\left(\frac{4569}{6250}\right)^2\frac65
=\frac{35028967}{97656250}<0.39<a.
\label{hmtfg:escalado:eq:21}
\end{equation}
Las órbitas de la hoja estable satisfacen \eqref{hmtfg:escalado:eq:19} a toda profundidad. Así, las iteraciones analíticas del operador corresponden a retornos reales de duplicación, con sus dominios y orientaciones conservados.


\subsection{Escalado de la cascada superestable local}
\label{hmtfg:escalado:8}

La conclusión métrica requiere dos cruces: el de la familia con la hoja estable, demostrado en sección~\ref{hmtfg:escalado:6}, y el de la variedad inestable con el blanco superestable. El segundo está probado por Eckmann–Wittwer para
\(\Sigma_2=\{\psi:\psi(1)=0\}\), correspondiente al ciclo superestable de período dos. La normalización \eqref{hmtfg:escalado:eq:6} coincide con la del operador real par empleado allí.

El teorema general de Collet–Eckmann–Lanford, sección 6 de la fuente citada, proposición 6.1 y teoremas 6.2–6.3, aplica a un mapa de Banach \(C^2\), un punto fijo con un único autovalor inestable \(\delta_{\mathrm F}>1\), una curva transversal a su hoja estable y un blanco transversal a su variedad inestable. Las preimágenes locales del blanco intersectan la curva en puntos que satisfacen
\begin{equation}
\delta_{\mathrm F}^j(\mu_j-\lambda_*)\longrightarrow C\ne0.
\label{hmtfg:escalado:eq:22}
\end{equation}
El índice \(j\) cuenta los retornos locales. Como la curva inicial aquí es \(R^4f_\lambda\), el período físico es \(2^{j+5}\) cuando el blanco se toma en \(\Sigma_2\); cualquier número fijo de pasos empleado para traer el blanco al entorno local se incorpora en un desplazamiento fijo del índice y en \(C\).

Las condiciones reales de sección~\ref{hmtfg:escalado:7} y la rama real de duplicación del blanco fijan el período exacto y el itinerario. La sucesión se selecciona por esta continuación local orientada. Tomando diferencias en \eqref{hmtfg:escalado:eq:22} se obtiene el límite de razones de \eqref{hmtfg:escalado:eq:4}.

Fuentes de los dos resultados utilizados: \href{https://link.springer.com/article/10.1007/BF01013368}{Eckmann–Wittwer, \emph{A complete proof of the Feigenbaum conjectures}}; \href{https://people.math.harvard.edu/~knill/history/lanford/papers/ColletEckmannLanford.pdf}{Collet–Eckmann–Lanford, sección 6}.

\subsubsection{Resto asintótico de potencia}

La analiticidad de nuestra familia permite precisar \eqref{hmtfg:escalado:eq:22}. En la construcción de coordenadas normales de Collet–Eckmann–Lanford se aplanan únicamente la variedad estable, la variedad inestable y la hipersuperficie blanco. Estas tres subvariedades admiten cambios locales \(C^2\) con derivadas acotadas. El corte suave empleado se realiza en la coordenada escalar inestable; la foliación completa se obtiene después mediante el cambio construido.

En esas coordenadas, la corrección se escribe como una serie de incrementos \(w_j\) cuya norma \(C^1\) satisface
\[
\|w_j\|_{C^1}\le C_1\rho^j,\qquad 0<\rho<1.
\]
La regla de la cadena para el mapa cortado, con derivadas de órdenes uno y dos uniformemente acotadas, proporciona
\[
\|w_j\|_{C^2}\le C_2B^j
\]
para algún \(B\ge1\). La norma ponderada de la construcción controla la norma \(C^1\) ordinaria en el dominio considerado. Por interpolación,
\[
[Dw_j]_\beta\le C_3
\left(\rho^{1-\beta}B^\beta\right)^j.
\]
Se elige
\begin{equation}
0<\beta<\min\left\{1,
\frac{-\log\rho}{\log B-\log\rho}\right\}.
\label{hmtfg:escalado:eq:23}
\end{equation}
La serie converge en \(C^{1,\beta}\). Sea \(z\) la coordenada normal resultante, normalizada para que las preimágenes del blanco satisfagan \(z=\delta_{\mathrm F}^{-j}\), absorbiendo el desplazamiento fijo de índice. Como nuestra curva es analítica y transversal,
\[
z(\Gamma(\lambda))=c(\lambda-\lambda_*)
+O(|\lambda-\lambda_*|^{1+\beta}),\qquad c\ne0.
\]
La inversión local da
\begin{equation}
\mu_j-\lambda_*=C\delta_{\mathrm F}^{-j}
+O\left(\delta_{\mathrm F}^{-(1+\beta)j}\right).
\label{hmtfg:escalado:eq:24}
\end{equation}
Esto demuestra el resto de \eqref{hmtfg:escalado:eq:4}, con \(\theta=\delta_{\mathrm F}^{-\beta}\). El enunciado afirma la existencia de las constantes del resto paramétrico. La tasa operatorial explícita \(0.83996\) de \eqref{hmtfg:escalado:eq:3} controla, por su parte, la convergencia de los mapas renormalizados.

\subsubsection{Conversión del resto en error de las razones}

Si \(\mu_j-\lambda_*=C\delta_{\mathrm F}^{-j}(1+e_j)\), con \(|e_j|\le M\theta^j\), se escribe
\[
\mu_{j-1}-\mu_j=C(\delta_{\mathrm F}-1)\delta_{\mathrm F}^{-j}(1+r_j),
\quad
r_j=\frac{\delta_{\mathrm F} e_{j-1}-e_j}{\delta_{\mathrm F}-1}.
\]
Con \(K=M(\delta_{\mathrm F}+\theta)/(\delta_{\mathrm F}-1)\), se tiene \(|r_j|\le K\theta^{j-1}\). Por tanto,
\begin{equation}
\boxed{
\left|\frac{\mu_{j-2}-\mu_{j-1}}{\mu_{j-1}-\mu_j}-\delta_{\mathrm F}\right|
\le\frac{\delta_{\mathrm F} K(1+\theta)\theta^{j-2}}
{1-K\theta^{j-1}}}
\label{hmtfg:escalado:eq:25}
\end{equation}
cuando \(K\theta^{j-1}<1\). Esta fórmula separa la precisión de evaluación de cada raíz del error de truncamiento de la cascada.


\subsection{Continuación certificada de las ocho raíces iniciales}
\label{hmtfg:escalado:9}

El certificado de continuación finita conserva la familia \eqref{hmtfg:escalado:eq:1} y el certificado anterior de ocho raíces. Para \(2\le n\le8\), estudia
\[
S_n(\lambda)=f_\lambda^{2^n}(0)
\]
desde la caja certificada de \(\lambda_{n-1}\) hasta \(1.874039\). La cobertura racional completa contiene exactamente la raíz heredada \(\lambda_{n-1}\) y la raíz nueva \(\lambda_n\). En particular, \(\lambda_n\) es la única raíz nueva en \((\lambda_{n-1},1.874039]\).

Las cajas se excluyen por signo del retorno o por una derivada local de signo certificado. Las cajas ancla conservan existencia y unicidad por signos extremos y monotonía local. Sus extremos cubren exactamente cada intervalo; se verifican las adyacencias. La evaluación usa los momentos \eqref{hmtfg:escalado:eq:8}, Horner de orden 32, colas geométricas exteriores y \(|u_0''|\le2\). La cola uniforme del potencial es menor que \(1.294\,10^{-58}\) en \(|x|\le3/2\).

El certificado final contiene 650 cajas procesadas y 332 hojas de cobertura. El primer nivel tiene la expresión exacta \(\lambda_1=1/u_0(1)\). Las aproximaciones siguientes se muestran únicamente como orientación; sus cajas completas se conservan en los certificados:

\begingroup\small
\begin{longtable}{p{0.465\linewidth}p{0.465\linewidth}}
\hline
Nivel & Parámetro superestable \\
\hline
1 & 1.345913990471832792210949095… \\
2 & 1.763379787846910127290794030… \\
3 & 1.850413796950136464530566778… \\
4 & 1.868979756606798125150574309… \\
5 & 1.872955034435659740004205128… \\
6 & 1.873806367467030119412262311… \\
7 & 1.873988695221365362307168780… \\
8 & 1.874027744154450672897007573… \\
\hline\end{longtable}
\endgroup
